\documentclass[10pt,a4paper]{amsart}
\usepackage[margin=19mm]{geometry}
\usepackage{amsmath,amssymb,mathtools}
\usepackage[scr=rsfs,cal=euler]{mathalfa}
\usepackage{xfrac}
\usepackage[unicode,hidelinks,bookmarks=false]{hyperref}
\newcommand{\ie}{i.e.\ }
\newcommand{\cf}{cf.\ }
\newcommand{\resp}{respectively\ }
\newcommand{\Subsection}{\subsection}
\providecommand{\nobreakdash}{\nobreak\hbox{-}}
\setlength{\emergencystretch}{2em}
\multlinegap0pt
\usepackage{xcolor}
\usepackage{amssymb}
\usepackage{mathtools}
\usepackage{bm}
\usepackage{dsfont}
\usepackage{xfrac}
\usepackage{algorithm}\usepackage{algpseudocode}
\usepackage{float}
\usepackage[shortlabels]{enumitem}
\usepackage{tikz}
\newtheorem{theorem}{Theorem}
\title{On the Capacity of Distinguishable Synthetic Identity Generation under Face Verification}
\author{Behrooz Razeghi}
\date{Source: arXiv:2604.10641v2 (CC BY 4.0)}
\begin{document}
\maketitle

\noindent\emph{Source and licence.} On the Capacity of Distinguishable Synthetic Identity Generation under Face Verification. Version arXiv:2604.10641v2 (CC BY 4.0), distributed by arXiv under the Creative Commons Attribution 4.0 International licence, http://creativecommons.org/licenses/by/4.0/, as recorded in the arXiv licence metadata for this exact version and shown on the abstract page https://arxiv.org/abs/2604.10641v2. Full licence notice: \texttt{licenses/arxiv-2604.10641-CC-BY-4.0.txt}.\par\smallskip

\noindent\textit{Editorial scope.} Counted unit: the article's theorem thm:full\_cap\_iff (source0.tex lines 2149--2185). The article's complete original proof of it is delivered with it (source0.tex lines 2187--2205, one \texttt{proof} environment of 19 lines). No mathematics is written, repaired or re-derived: the statement and the proof are the article's own lines, in the article's own notation, and no retained line is substituted or re-flowed.  No further source-local context is needed: every cross-reference of every retained line resolves inside the two delivered spans.  Nothing was omitted from any retained span, so the package carries no omission record.  No retained line contains a citation, so the wrapper carries no bibliography and no bibliography entry.  No retained line contains a figure environment, an image-inclusion command or drawing code, so no image is delivered and the package declares no asset dependency.  Editorial formatting only, in addition: an AMS \texttt{amsart} preamble that restates the article's own declarations and macros used by the retained lines, namely the environments \texttt{theorem} on separate counters, together with the standard packages those lines need.  The retained environments print in this document as Theorem 1 in order of appearance, whereas the article prints the same items with the numbering of its own full document.  The complete native source of the article as distributed in the arXiv e-print for this exact version is bundled unchanged as \texttt{sources/198285/source0.tex}.
\begin{theorem}[Full-cap zero-error conditions]
\label{thm:full_cap_iff}
%
Fix $\rho\in[0,\pi/2)$ and $\tau\in[0,1)$, and suppose $\{P_i\}_{i=1}^M$ is a family such that $P_i$ is $\rho$-full-cap-supported at $u_i\in\mathbb S^{D-1}$ for every $i\in[M]$.
%
\begin{enumerate}
%
\item If the family is $(\tau,0,0)$-admissible, then
%
\begin{align}
2\rho &\le \arccos(\tau),
\label{eq:full_cap_intra_iff}\\
\angle(u_i,u_j) &\ge \arccos(\tau)+2\rho, \qquad i\ne j.
\label{eq:full_cap_inter_iff}
\end{align}
%
\item If
%
\begin{equation}
2\rho<\arccos(\tau)
\label{eq:full_cap_intra_suff}
\end{equation}
%
and \eqref{eq:full_cap_inter_iff} holds, then the family is $(\tau,0,0)$-admissible.
%
\end{enumerate}
%
At the boundary $2\rho=\arccos(\tau)$, every within-cap pair has cosine score at least $\tau$, but zero-error genuine acceptance additionally requires
%
\begin{equation}
(P_i\otimes P_i) \left( \left\{ (e,e')\in \mathbb S^{D-1}\times\mathbb S^{D-1} \,\middle|\,
\langle e,e'\rangle=\tau \right\} \right) =0
\end{equation}
%
for every $i\in[M]$. Thus support alone does not determine the zero-error genuine condition at the boundary.
%
\end{theorem}

\begin{proof}
%
For necessity of \eqref{eq:full_cap_intra_iff}, suppose $2\rho>\arccos(\tau)$. The cap $\mathrm{Cap}(u_i,\rho)$ contains two points whose angular separation is $2\rho$ and whose cosine score is therefore strictly smaller than $\tau$. By continuity of the cosine score, there exist relative neighborhoods of these points within the cap such that every pair formed by one point from each neighborhood has score strictly smaller than $\tau$. Full support assigns positive $P_i$-probability to each neighborhood. Since the two genuine samples are independent, the event that they lie in the respective neighborhoods has positive probability. The genuine rejection probability is therefore positive, contradicting $(\tau,0,0)$-admissibility. Hence $2\rho\le\arccos(\tau)$.

For necessity of \eqref{eq:full_cap_inter_iff}, suppose for some $i\ne j$ that
%
\begin{equation}
\angle(u_i,u_j) < \arccos(\tau)+2\rho.
\end{equation}
%
Then $\mathrm{Cap}(u_i,\rho)$ and $\mathrm{Cap}(u_j,\rho)$ contain points whose angular separation is strictly smaller than $\arccos(\tau)$. Their cosine score is therefore strictly larger than $\tau$. By continuity, there exist relative neighborhoods of these points within the respective caps such that every corresponding cross-cap pair has score strictly larger than $\tau$. Full support assigns positive probability to both neighborhoods. Independence of the cross-identity samples therefore gives a positive-probability impostor match event, contradicting $(\tau,0,0)$-admissibility. Hence \eqref{eq:full_cap_inter_iff} is necessary.

For sufficiency, \eqref{eq:full_cap_intra_suff} implies that every pair of points in the same radius-$\rho$ cap has angular separation at most $2\rho<\arccos(\tau)$ and therefore cosine score strictly larger than $\tau$. Thus the genuine constraint holds with probability one. Under \eqref{eq:full_cap_inter_iff}, every pair of points from distinct caps has angular separation at least
\[
\angle(u_i,u_j)-2\rho \ge \arccos(\tau)
\]
and therefore cosine score at most $\tau$. Thus the impostor non-match constraint also holds with probability one. The family is consequently $(\tau,0,0)$-admissible.
%
\end{proof}

\end{document}
